\documentclass[12pt]{article}
\usepackage[utf8]{inputenc}
\usepackage[T1]{fontenc}
\usepackage{hyperref}
\title{Majera Language Resources for K-12 Curriculum Development}
\author{Compiled by Claude Sonnet 5.5 (claude-sonnet-5-5) for Andr\'as Kornai}
\date{}
\begin{document}
\maketitle

\section{Introduction}
Majera (ISO 639-3 xmj; Glottocode maje1243), also called Midah, Mida'a or Daa, is a Chadic (Biu-Mandara) language of the Afro-Asiatic family. It is spoken around Majera in the Zina arrondissement, Logone-et-Chari, Far North Region of Cameroon, and in Chad. Estimates of speaker numbers vary widely: Wikipedia cites 500 (2004), Joshua Project gives 900 in Cameroon and 1,600 in Chad. The language is considered endangered. Search turned up only general reference pages.

\section{Primary Educational Resources}
No primers, textbooks or school programmes were found.

\section{Linguistic and Reference Materials}
\begin{itemize}
\item Wikipedia, ``Majera language'': \url{https://en.wikipedia.org/wiki/Majera_language}
\item Endangered Languages Project entry: \url{https://endangeredlanguages.com/lang/2305?hl=de}
\item Lexvo entry for xmj: \url{https://finto.fi/lexvo/en/page/xmj}
\end{itemize}

\section{Religious and Community Resources}
\begin{itemize}
\item Joshua Project, Majera in Cameroon: \url{https://www.joshuaproject.net/index.php/people_groups/13230/CM}
\item Joshua Project, Majera in Chad: \url{https://www.joshuaproject.net/index.php/people_groups/13230/CD}
\end{itemize}

\section{Digital and Audio Resources}
None found in the search results.

\section{Curriculum Development Context}
The language is not known to be taught in schools. No orthography or teaching materials were found, so a curriculum would require documentation work with the community first. Because speakers are few and shifting to other languages, even basic documentation would be valuable.
\end{document}